\documentclass[12pt,a4paper]{article}
\usepackage[T1]{fontenc}
\usepackage{amsmath,amssymb,mathrsfs,tikz,times,pifont}
\usepackage{enumitem}
\usepackage{multicol}
\usepackage{lmodern}
\usetikzlibrary{trees}
\newcommand\circitem[1]{%
\tikz[baseline=(char.base)]{
\node[circle,draw=gray, fill=red!55,
minimum size=1.2em,inner sep=0] (char) {#1};}}
\newcommand\boxitem[1]{%
\tikz[baseline=(char.base)]{
\node[fill=cyan,
minimum size=1.2em,inner sep=0] (char) {#1};}}
\setlist[enumerate,1]{label=\protect\circitem{\arabic*}}
\setlist[enumerate,2]{label=\protect\boxitem{\alph*}}
%%%::::::by chnini ameur :::::::%%%
\everymath{\displaystyle}
\usepackage[left=1cm,right=1cm,top=1cm,bottom=1.7cm]{geometry}
\usepackage[colorlinks=true, linkcolor=blue, urlcolor=blue, citecolor=blue]{hyperref}
\usepackage{array,multirow}
\usepackage[most]{tcolorbox}
\usepackage{varwidth}
\usepackage{float} %pour utiliser l'option [H] qui force l'image à apparaître exactement à l'endroit où elle est placée dans le code.
\tcbuselibrary{skins,hooks}
\usetikzlibrary{patterns}
%%%::::::by chnini ameur :::::::%%%
\newtcolorbox{exa}[2][]{enhanced,breakable,before skip=2mm,after skip=5mm,
colback=yellow!20!white,colframe=black!20!blue,boxrule=0.5mm,
attach boxed title to top left ={xshift=0.6cm,yshift*=1mm-\tcboxedtitleheight},
fonttitle=\bfseries,
title={#2},#1,
% varwidth boxed title*=-3cm,
boxed title style={frame code={
\path[fill=tcbcolback!30!black]
([yshift=-1mm,xshift=-1mm]frame.north west)
arc[start angle=0,end angle=180,radius=1mm]
([yshift=-1mm,xshift=1mm]frame.north east)
arc[start angle=180,end angle=0,radius=1mm];
\path[left color=tcbcolback!60!black,right color = tcbcolback!60!black,
middle color = tcbcolback!80!black]
([xshift=-2mm]frame.north west) -- ([xshift=2mm]frame.north east)
[rounded corners=1mm]-- ([xshift=1mm,yshift=-1mm]frame.north east)
-- (frame.south east) -- (frame.south west)
-- ([xshift=-1mm,yshift=-1mm]frame.north west)
[sharp corners]-- cycle;
},interior engine=empty,
},interior style={top color=yellow!5}}
%%%%%%%%%%%%%%%%%%%%%%%
\usepackage{fancyhdr}
\usepackage{eso-pic}         % Pour ajouter des éléments en arrière-plan
% Commande pour ajouter du texte en arrière-plan
\usepackage{tkz-tab}
\AddToShipoutPicture{
    \AtTextCenter{%
        \makebox[0pt]{\rotatebox{80}{\textcolor[gray]{0.7}{\fontsize{5cm}{5cm}\selectfont PGB}}}
    }
}
\usepackage{lastpage}
\fancyhf{}
\pagestyle{fancy}
\renewcommand{\footrulewidth}{1pt}
\renewcommand{\headrulewidth}{0pt}
\renewcommand{\footruleskip}{10pt}
\fancyfoot[R]{
\color{blue}\ding{45}\ \textbf{2025}
}
\fancyfoot[L]{
\color{blue}\ding{45}\ \textbf{Prof:M. BA}
}
\cfoot{\bf
\thepage /
\pageref{LastPage}}
% Création du compteur pour les exercices
\newcounter{exercice}
\renewcommand{\theexercice}{\arabic{exercice}}  % Définit l'affichage du compteur en chiffres arabes

% Définir la commande \exo
\newcommand{\exo}{\refstepcounter{exercice}\textbf{Exercice \theexercice} }

\begin{document}
\renewcommand{\arraystretch}{1.5}
\renewcommand{\arrayrulewidth}{1.2pt}
\begin{tikzpicture}[overlay,remember picture]
    \node[draw=blue,line width=1.2pt,fill=purple,text=blue,inner sep=3mm,rounded corners,pattern=dots]at ([yshift=-2.5cm]current page.north) {\begingroup\setlength{\fboxsep}{0pt}\colorbox{white}{\begin{tabular}{|*1{>{\centering \arraybackslash}p{0.28\textwidth}} |*2{>{\centering \arraybackslash}p{0.2\textwidth}|} *1{>{\centering \arraybackslash}p{0.19\textwidth}|} }
                \hline
                \multicolumn{3}{|c|}{$\diamond$$\diamond$$\diamond$\ \textbf{Lycée de Dindéfélo}\ $\diamond$$\diamond$$\diamond$ } & \textbf{A.S. : 2025/2026}                                              \\ \hline
                \textbf{Matière: Mathématiques}                                                                                    & \textbf{Niveau : 2}\textbf{ndS} & \textbf{Date: 29/12/2025} & \textbf{} \\ \hline
                \multicolumn{4}{|c|}{\parbox[c]{10cm}{\begin{center}
                                                                  \textbf{{\Large\sffamily Td Calcul vectoriel}}
                                                              \end{center}}}                                                                                                        \\ \hline
            \end{tabular}}\endgroup};
\end{tikzpicture}
\vspace{3cm}

\fbox{\textbf{\exo}}


\textbf{1. Expression de $\overrightarrow{IJ}$ en fonction de $\overrightarrow{AB}$ et $\overrightarrow{AC}$}
En utilisant la relation de Chasles :
\[ \overrightarrow{IJ} = \overrightarrow{IA} + \overrightarrow{AC} + \overrightarrow{CJ} \]
\begin{itemize}
    \item Comme $I$ est le milieu de $[AB]$, on a $\overrightarrow{IA} = \frac{1}{2}\overrightarrow{BA} = -\frac{1}{2}\overrightarrow{AB}$.
    \item La relation de l'énoncé donne $\overrightarrow{CJ} = \frac{1}{4}\overrightarrow{AC}$.
\end{itemize}
D'où :
\[ \overrightarrow{IJ} = -\frac{1}{2}\overrightarrow{AB} + \overrightarrow{AC} + \frac{1}{4}\overrightarrow{AC} \]
\[ \overrightarrow{IJ} = -\frac{1}{2}\overrightarrow{AB} + \frac{5}{4}\overrightarrow{AC} \]

\textbf{2. Expression de $\overrightarrow{IK}$ en fonction de $\overrightarrow{AB}$ et $\overrightarrow{AC}$}
De même :
\[ \overrightarrow{IK} = \overrightarrow{IB} + \overrightarrow{BC} + \overrightarrow{CK} \]
\begin{itemize}
    \item $\overrightarrow{IB} = \frac{1}{2}\overrightarrow{AB}$.
    \item $\overrightarrow{BC} = \overrightarrow{BA} + \overrightarrow{AC} = -\overrightarrow{AB} + \overrightarrow{AC}$.
    \item $\overrightarrow{CK} = \frac{1}{6}\overrightarrow{CB} = \frac{1}{6}(\overrightarrow{CA} + \overrightarrow{AB}) = -\frac{1}{6}\overrightarrow{AC} + \frac{1}{6}\overrightarrow{AB}$.
\end{itemize}
En sommant :
\[ \overrightarrow{IK} = \left( \frac{1}{2} - 1 + \frac{1}{6} \right)\overrightarrow{AB} + \left( 1 - \frac{1}{6} \right)\overrightarrow{AC} \]
\[ \overrightarrow{IK} = -\frac{1}{3}\overrightarrow{AB} + \frac{5}{6}\overrightarrow{AC} \]

\textbf{3. Conclusion sur l'alignement}
Cherchons s'il existe un réel $k$ tel que $\overrightarrow{IK} = k \cdot \overrightarrow{IJ}$. \\
Comparons les coefficients :
\[ \frac{-\frac{1}{3}}{-\frac{1}{2}} = \frac{2}{3} \quad \text{et} \quad \frac{\frac{5}{6}}{\frac{5}{4}} = \frac{5}{6} \times \frac{4}{5} = \frac{4}{6} = \frac{2}{3} \]
On a donc $\overrightarrow{IK} = \frac{2}{3}\overrightarrow{IJ}$. Les vecteurs $\overrightarrow{IJ}$ et $\overrightarrow{IK}$ sont colinéaires et ont le point $I$ en commun. \\
\textbf{Les points $I$, $J$ et $K$ sont donc alignés.}


\end{document}